% !PRINTED 371
\weberpage{371}

\begin{center}
§. 120.

\gesperrt{Aequivalente Zahlen.}
\end{center}

Die Kettenbrüche führen uns auf die Betrachtung einer
besonderen Art von linearen Substitutionen, die in der Algebra
und Zahlentheorie überhaupt eine grosse Bedeutung haben, und
die wir etwas näher betrachten wollen.

Sind $x$ und $y$ zwei Zahlen oder auch veränderliche Grössen,
die in der Abhängigkeit von einander stehen
\weq{1}{y = \frac{\alpha x + \beta}{\gamma x + \delta},}

worin $\alpha, \beta, \gamma, \delta$ ganze Zahlen sind, die der Bedingung
\weq{2}{\alpha \delta - \beta \gamma = \varepsilon = \pm 1}

genügen, so nennen wir $x$ und $y$ mit einander \gesperrt{äquivalent}$^{1)}$.
Wir nennen sie \gesperrt{eigentlich} oder \gesperrt{uneigentlich äquivalent},
je nachdem $\varepsilon = +1$ oder $\varepsilon = -1$ ist. Diese Beziehung ist
eine gegenseitige, denn aus (1) folgt
\weq{3}{x = \frac{\delta y - \beta}{-\gamma y + \alpha}.}

\gesperrt{Wenn zwei Grössen mit einer dritten äquivalent
sind, so sind sie auch mit einander äquivalent.}

Denn ist
\weq{4}{x = \frac{\alpha' z + \beta'}{\gamma' z + \delta'},}

so folgt aus (1)
\weq{5}{y = \frac{\alpha'' z + \beta''}{\gamma'' z + \delta''},}

wenn
\weq{6}{\begin{aligned}
\alpha'' &= \alpha \alpha' + \beta \gamma', \quad \beta'' = \alpha \beta' + \beta \delta', \\
\gamma'' &= \gamma \alpha' + \delta \gamma', \quad \delta'' = \gamma \beta' + \delta \delta',
\end{aligned}}
\weq{7}{\alpha'' \delta'' - \beta'' \gamma'' = (\alpha \delta - \beta \gamma)(\alpha' \delta' - \beta' \gamma'),}

oder
\weq{8}{\varepsilon'' = \varepsilon \varepsilon'.}

\begingroup
\renewcommand{\thefootnote}{1)}
\footnotetext[1]{Vgl. \gesperrt{Dedekind}, Schreiben an Herrn \gesperrt{Borchardt} über die Theorie
der elliptischen Modulfunctionen. Crelle's Journal, Bd. 83 (1877).}
\endgroup
